\chapter{Motivation}%
\label{s_motivation}

\section{Importance of Scalable Path Planning in U-Space}%
\label{s_motivation}

Urban Air Mobility (UAM) concepts rely on safe and efficient trajectory planning 
within highly dynamic and dense low-altitude airspace. The European U-space 
framework envisions a digital and service-oriented architecture to enable such 
operations \cite{eu_commission_2017}. 

A key requirement for this environment is scalability: thousands of trajectory 
requests may need to be processed in real time, while ensuring separation 
minima and maintaining safe operations. Previous research has shown that 
trajectory-based operations require computationally efficient frameworks to 
handle large volumes of requests without degrading performance 
\cite{kopardekar2016utm, sesar_conops_2019}. 

\section{Challenges in Real-Time Multi-UAV Coordination}%
\label{s_contribution}

Path planning interacts with several other services in the UAM ecosystem. 
The global planner provides a feasible route across the corridor network, 
but this must be complemented by flight management software capable of 
synthesizing flyable 4D trajectories \cite{morscheck2021cfms}, conflict 
detection modules such as NDMap \cite{NDMap}, and path verification 
services \cite{sesar_conops_2019}. 

If these services are tightly coupled in a monolithic architecture, the system 
risks reduced interoperability, higher technical debt, and limited extensibility. 
Given the diversity of programming tools and modeling approaches used in these 
modules, there is a strong need for a framework that allows independent 
implementations to interact seamlessly.

\section{Role of Microservices in UTM Systems}%
\label{s_structure}

A microservice-based design allows path planning to be implemented as an 
independent service, decoupled from trajectory generation, verification, and 
conflict detection. This promotes modularity, extensibility, and ease of 
maintenance \cite{dragoni2017microservices, matlekovic2020uav}. 

Such an approach is particularly relevant in the context of U-space, where 
real-time performance and scalability are critical. Evaluating this framework 
in terms of scalability and latency is essential to assess its viability for 
future UAM applications \cite{kopardekar2016utm}. 

Finally, the usability of these systems also depends on visualization and 
stakeholder access. Different actors (e.g., operators, regulators, or 
airspace managers) may require different levels of visibility into the 
system. Determining which information should be accessible to which user 
group remains an open question of both technical and regulatory importance 
\cite{sesar_master_plan_2025}.
